\section{The work of Ki and the Thorin dictionary}\label{ki:sec:dict}

\subsection{The strands of the work}

The work of Ki and his coauthors falls into four strands that matter here. The first concerns the heat flow on Fourier transforms of rapidly decaying kernels: Ki and Kim extend a result of Newman on the zeros of Fourier transforms \cite{KiKim02} and answer a question of de Bruijn \cite{KiKim03}, and Ki, Kim and Lee prove that for every positive time the flowed function has all but finitely many of its zeros real and simple \cite{KKL09}. % UNCHECKED: KiKim02 generalises which result of Newman; statements of KiKim02 and KiKim03 not checked against the papers
The second concerns repeated differentiation: Ki and Kim settle the Fourier P\'olya conjecture \cite{KiKim00}, and Ki studies the Riemann $\Xi$ function under repeated differentiation \cite{Ki06b}. The third concerns approximate zeta functions, which keep the functional equation of a zeta function and lose its Euler product: approximations of the Epstein zeta function \cite{Ki05}, modified Epstein zeta functions \cite{Ki06a}, approximations of the Ramanujan $\Xi$ function \cite{Ki08a}, Weng's zeta functions \cite{Ki10} and, with Komori and Suzuki, Weng zeta functions for Chevalley groups \cite{KKS15}. The fourth concerns the zeros of $\zeta'$ \cite{Ki08b,FarmerKi12,GonekKi18} and is taken up in Section~\ref{ki:sec:open}.

\subsection{Zeros as Thorin atoms}

The P\'olya kernel $\Phi$ and the function $\Xi(\theta)=\xi(\tfrac12+i\theta)=\int_{\R}\Phi(u)\cos(\theta u)\,du$ are those of \eqref{eq:ch2:Phi}, and $X$ is the P\'olya variable with density $\Phi/\Xi(0)$, so that $\Xi(\theta)/\Xi(0)=\E[e^{i\theta X}]$. With $s=-\theta^2$ this reads $\Xi(i\sqrt s)/\Xi(0)=\E[\cosh(\sqrt s\,X)]=G(s)$, the function of Section~\ref{ch:clock}, because $\xi(\tfrac12-\sqrt s)=\xi(\tfrac12+\sqrt s)$. Ki's results concern real entire functions of this kind, and the following lemma carries the Thorin reading over to all of them.

\begin{lemma}\label{ki:lem:dict}
\begin{enumerate}
\item[(a)] Let $F$ be an even entire function, real on $\R$, of order less than two, with $F(0)\ne0$. Let $\pm w_k$ be its zeros, with multiplicity (one sign chosen for each pair, and $w_k\ne0$), and put $\lambda_k=w_k^2$. Then $s\mapsto F(i\sqrt s)$ is entire of order less than one and
\[
\frac{F(0)}{F(i\sqrt s)}=\prod_k\Bigl(1+\frac{s}{\lambda_k}\Bigr)^{-1},
\]
locally uniformly off the points $-\lambda_k$. If every $w_k$ is real, this is the Laplace transform of $T=\sum_k\lambda_k^{-1}e_k$, with $e_k$ independent standard exponential variables: a GGC with zero drift and discrete Thorin measure $\sum_k\delta_{\lambda_k}$. If some $w_k$ is not real, then $\lambda_k\notin(0,\infty)$, and $F(0)/F(i\sqrt s)$ is not the Laplace transform of any GGC.
\item[(b)] Let $K$ be even and positive, with $K(u)=O(\exp(-|u|^p))$ for some $p>2$, and let $F(z)=\int_{\R}K(u)e^{izu}\,du$. Then $F$ satisfies the hypotheses of (a).
\end{enumerate}
\end{lemma}

\begin{proof}
(a) The function $f(s)=F(i\sqrt s)$ is entire because $F$ is even, and it has the same order as $F$ divided by two, so its order is less than one. Hadamard's theorem gives $f(s)=f(0)\prod_k(1+s/\lambda_k)$ with convergent product, because an entire function of order less than one has genus zero. The product is the Laplace transform of $T$ when all $\lambda_k>0$, since $\E e^{-se_k/\lambda_k}=(1+s/\lambda_k)^{-1}$. A zero $w_k\notin\R$ has $w_k^2\notin(0,\infty)$. A GGC transform $\exp\{-\int\log(1+s/t)\,\mu(dt)\}$ with $\mu\ge0$ on $(0,\infty)$ is analytic on $\C\setminus(-\infty,0]$, while $1/f$ has a pole at $-\lambda_k\notin(-\infty,0]$, of the order of the multiplicity of the zero, which no other factor cancels.

(b) $F(0)=\int K>0$. The bound on $K$ gives $|F(x+iy)|\le C\int e^{|y||u|-|u|^p}du\le C'\exp(c_p|y|^{p/(p-1)})$, so $F$ has order at most $p/(p-1)<2$, and $F$ is even and real.
\end{proof}

For $F=\Xi$ the zeros $\pm\theta$ are $\pm\gamma$ with $\gamma>0$, and Lemma~\ref{ki:lem:dict}(a) gives $\xi(\tfrac12)/\xi(\tfrac12+\sqrt s)=\prod_{\gamma>0}(1+s/\gamma^2)^{-1}$, the clock of Section~\ref{ch:clock}. RH is the statement that its Thorin measure $\sum_{\gamma>0}\delta_{\gamma^2}$ is a positive measure on $(0,\infty)$. Special cases of the lemma, and the pole argument in its proof, appear in Propositions~\ref{prop:rigid}, \ref{rz:prop:newmanthorin} and \ref{pk:prop:filter}; the general form is needed for the approximations of Section~\ref{ki:sec:approx}. The atoms of a hypothetical counterexample are confined by a simple map.

\begin{lemma}[Strips become confocal parabolas]\label{ki:lem:parab}
For $\Delta>0$ the map $w\mapsto w^2$ sends the strip $\{|\im w|\le\Delta\}$ onto the closed region to the right of the parabola $\mathcal P_\Delta=\{x+iy:\ x=y^2/(4\Delta^2)-\Delta^2\}$. Every $\mathcal P_\Delta$ has its focus at the origin, the parabolas are nested, and they collapse onto $[0,\infty)$ as $\Delta\downarrow0$.
\end{lemma}

\begin{proof}
The line $w=a+ib$, $a\in\R$, $0<b\le\Delta$, goes to $w^2=(a^2-b^2)+2iab$, that is, to the parabola $x=y^2/(4b^2)-b^2$. For fixed $y\ne0$ the right side decreases from $+\infty$ to $y^2/(4\Delta^2)-\Delta^2$ as $b$ increases from $0$ to $\Delta$, so the images of the lines with $0<b\le\Delta$ fill the points with $x\ge y^2/(4\Delta^2)-\Delta^2$, $y\ne0$. The real axis ($b=0$) goes to $[0,\infty)$, and the points $x\in[-\Delta^2,0)$ on the axis $y=0$ are the images of $w=i\sqrt{-x}$. The lines with $b\le\Delta$ lie in the strip and the same computation shows that the image has no point left of $\mathcal P_\Delta$. The focus of $x=y^2/(4\Delta^2)-\Delta^2$ is at $x=-\Delta^2+\Delta^2=0$. For $\Delta'<\Delta$ the parabola $\mathcal P_{\Delta'}$ lies to the right of $\mathcal P_\Delta$ at every height $y\ne0$, and for fixed $y\ne0$ its abscissa tends to $+\infty$ as $\Delta'\downarrow0$.
\end{proof}

\section{The heat flow, the barrier and the Gaussian tilt}\label{ki:sec:flow}

For real $t$ let $\Xi_t$ be the transform of the damped kernel $e^{tu^2/4}\Phi(u)$ of Section~\ref{sec:biane:flow}, so that $\Xi_t(\theta)=\int_{\R}\Phi(u)e^{tu^2/4}\cos(\theta u)\,du$ and $\Xi_0=\Xi$. With $D=d/d\theta$ we have $D^2\cos(\theta u)=-u^2\cos(\theta u)$, hence $\Xi_t=e^{-(t/4)D^2}\Xi$: the backward heat operator is the Gaussian tilt of the P\'olya density, and $\Xi_t(\theta)=8H_t(2\theta)$ with $H_t$ the function of Section~\ref{sec:ch21:flow}, so that the constants of that section carry over without change. The zeros of $\Xi$ lie in $|\im\theta|<\tfrac12$ and the zeros of $\Xi_t$ in $(\im\theta)^2\le\max(\tfrac14-\tfrac t2,0)$ by de Bruijn's theorem \cite[Theorem~13]{deBruijn} (Theorem~\ref{tf:thm:heat}(c) with $r=1$, since the zeros of $H_0$ lie in $|\im z|\le1$); in particular $\Xi_t$ has only real zeros for $t\ge\tfrac12$. Newman showed that there is a constant $\Lambda$ such that $\Xi_t$ has only real zeros exactly for $t\ge\Lambda$ \cite{Newman1976}, and Rodgers and Tao proved $\Lambda\ge0$ \cite{RodgersTao}; RH is $\Lambda\le0$. By Lemma~\ref{ki:lem:dict}(b), applied to $K=\Phi e^{tu^2/4}$, the atoms of the Thorin measure of $1/\Xi_t$ are the squares of the zeros of $\Xi_t$ in the upper half plane and on the positive axis.

\subsection{The theorem of Ki, Kim and Lee in Thorin form}\label{ki:sec:kkl}

Ki, Kim and Lee proved that for every $t>0$ the function $\Xi_t$ has infinitely many zeros, all but finitely many of them real and simple \cite{KKL09}; this gives $\Lambda<\tfrac12$ and no height beyond which the zeros are real (Remark~\ref{mw:rem:flowcert} records the same statement).
Lemma~\ref{ki:lem:dict} turns it into a statement on the Thorin measure.

\begin{proposition}[Eventually Thorin]\label{ki:prop:eventual}
For every $t>0$ there are a finite, conjugation invariant multiset $E_t\subset\C\setminus(0,\infty)$ and a GGC Laplace transform $L_t$, with discrete Thorin measure on $(0,\infty)$ whose atoms are simple beyond a finite level, such that
\[
\frac{\Xi_t(0)}{\Xi_t(i\sqrt s)}=L_t(s)\prod_{\eta\in E_t}\Bigl(1+\frac{s}{\eta}\Bigr)^{-1}.
\]
$\Xi_t$ has only real zeros exactly when $E_t=\emptyset$.
\end{proposition}

\begin{proof}
The kernel $\Phi e^{tu^2/4}$ is even and positive and decays like $\exp(-\pi e^{2u}/2)$, so Lemma~\ref{ki:lem:dict}(b) applies to it for every $p$. Split the zeros of $\Xi_t$ into the real ones and the finitely many nonreal ones, which by the theorem of Ki, Kim and Lee form a conjugation invariant set. The nonreal zeros $w$ give the multiset $E_t=\{w^2\}$, which misses $(0,\infty)$, and the real zeros give $L_t$ by Lemma~\ref{ki:lem:dict}(a). Only finitely many real zeros are multiple.
\end{proof}

For $t>0$ the tail of the Thorin measure of the reciprocal is therefore positive, and only a finite rational factor has to be settled by computation. The case $t=0$ of the identity would be the clock of $\xi$ with $E_0$ the multiset of squared off line zeros, and RH is the statement that $E_0$ is empty. The barrier of the Polymath project is a separate device, for the zeros that enter the strip from large heights (Chapter~\ref{chap:phase}); the Polymath project proved $\Lambda\le0.22$ \cite{Polymath} with an effective version of the theorem and a barrier, and Platt and Trudgian improved this to $\Lambda\le0.2$ \cite{PT21}.

De Bruijn's strip theorem says where $E_t$ can be; Ki and Kim \cite{KiKim03} proved the form of it that allows finitely many exceptional zeros: if for every $\varepsilon>0$ all but finitely many zeros of a real entire function $f$ of genus at most one lie in $|\im z|\le\Delta+\varepsilon$, then for every $\varepsilon>0$ all but finitely many zeros of $e^{-\lambda D^2}f$ lie in $|\im z|\le\sqrt{\max(\Delta^2-2\lambda,0)}+\varepsilon$. % UNCHECKED: KiKim03 statement as in the chapter source
That statement is about all but finitely many zeros and does not locate the finitely many exceptions, so by itself it adds nothing to Proposition~\ref{ki:prop:eventual}. What does locate $E_t$ is de Bruijn's theorem for all zeros.

\begin{corollary}[Shrinking parabolas]\label{ki:cor:parab}
For $0<t<\tfrac12$ every Thorin atom of $\Xi_t$, in particular every $\eta\in E_t$, lies in the closed region to the right of $\mathcal P_{\Delta(t)}$, where $\Delta(t)^2=\tfrac14-\tfrac t2$. The parabolas share the focus $0$ and collapse onto $[0,\infty)$ at $t=\tfrac12$, where de Bruijn's theorem takes over.
\end{corollary}

\begin{proof}
By Theorem~\ref{tf:thm:heat}(c) with $r=1$, all zeros of $\Xi_t$ lie in $(\im\theta)^2\le\tfrac14-\tfrac t2=\Delta(t)^2$, and an atom is the square of a zero. Lemma~\ref{ki:lem:parab} gives the region.
\end{proof}

\subsection{The Gaussian tilt obstruction}\label{ki:sec:obstruction}

Read through Lemma~\ref{ki:lem:dict}, the theorem of Rodgers and Tao says that for every $t<0$ the Thorin measure of $1/\Xi_t$ has atoms off $(0,\infty)$ (Proposition~\ref{ki:prop:obstruction}). This is a filter on proof strategies.

\begin{proposition}[No slack in the Gaussian direction]\label{ki:prop:obstruction}
Let $\mathcal K$ be a class of positive even kernels closed under $K\mapsto Ke^{-bu^2}$ for every $b>0$. If every $K\in\mathcal K$ has a Fourier transform with only real zeros, then $\Phi\notin\mathcal K$.
\end{proposition}

\begin{proof}
If $\Phi\in\mathcal K$, then $\Phi e^{-bu^2}\in\mathcal K$ for every $b>0$, so $\Xi_t$ has only real zeros for every $t<0$, that is, $\Lambda<0$, contradicting $\Lambda\ge0$ \cite{RodgersTao}.
\end{proof}

The statement is Proposition~\ref{prop:biane:damping}(a) and Proposition~\ref{pk:prop:filter} in other words, and the remark that follows is what the proposition adds. Multiplication by $e^{-bu^2}$ adds $-2b$ to $(\log K)''$ and the linear function $-bv$ to $\log K(\sqrt v)$. Every condition that bounds $(\log K)''$ from above is therefore preserved, including log concavity of $K(u)$ and of $K(\sqrt v)$. No such condition can imply RH together with any further hypotheses that the tilt also preserves. For log concavity this is Proposition~\ref{pk:prop:filter}(ii); see also \cite{Newton26}. The same holds for any form of the sine squared lemma (eq.~\eqref{eq:sin2}) whose positivity hypothesis on the L\'evy measure survives the corresponding tilt. A Thorin argument for RH must therefore use a property that $\Phi$ has and that $\Phi e^{-bu^2}$ loses for every $b>0$.

\begin{remark}[Newman's class]
Newman \cite{Newman74} characterised the Lee and Yang single spin laws as those whose Fourier transform is in the Laguerre P\'olya class (Proposition~\ref{rz:prop:newmanthorin}); the class is preserved by $e^{bt^2}$, $b>0$. The law with density $\Phi/\Xi(0)$ belongs to it if and only if RH holds, and by Rodgers and Tao $\Phi e^{-bu^2}$ does not belong to it for any $b>0$, so under RH the P\'olya law lies on its boundary.
\end{remark}
